% Section 12: verification methodology.
\section{Verification methodology}\label{sec:verification}

A solver is useful only if its ``optimal'' can be trusted. \nirnay{} is checked at four levels.

\subsection{Independent references, in separate processes}
The harnesses \file{bench/run\_lp.py}, \file{run\_mip.py} and \file{run\_qp.py} run every instance
in a fresh process with a wall-clock limit, so one pathological instance cannot stall or corrupt a
sweep. Each process first solves a tiny instance, so that loading the Numba-compiled code is not
charged to the first real instance, and then times only the solve call (file reading excluded).
The same file is then solved by a reference solver:
\begin{itemize}
  \item \textbf{LP and MILP:} HiGHS~\verHighs~\cite{huangfu2018} with the same time limit, reading its
        own copy of the file;
  \item \textbf{QP:} Clarabel~\verClarabel~\cite{goulart2024}, an interior-point conic solver. HiGHS's
        active-set QP solver times out on the larger Maros--Mészáros problems. The model Clarabel sees
        is read by \emph{HiGHS's} MPS parser, not ours, so a reader bug in \nirnay{} cannot hide by
        affecting both sides. The failure analysis of the first QP sweep found two defects in this
        reference, and the harness was corrected for the second sweep. First, at its default
        tolerances Clarabel's objective is about $10^{-6}$ off on several problems where \nirnay{}
        matches the published optimum, so the reference now runs at $10^{-11}$. Second, a comparison of
        the two parsers (\file{tools/qps\_compare.py}) showed that HiGHS~\verHighs{} misreads the RHS
        section of \code{DPKLO1} and drops coefficients below $10^{-9}$ in \code{HUESTIS},
        \code{HUES-MOD} and \code{KSIP}. On files where the two parsers disagree, the reference now
        solves \nirnay{}'s reading and says so in its status.
\end{itemize}
Neither reference is imported by the \code{nirnay} package. \file{bench/comparators.py}, which drives
the other solvers of \cref{sec:competitors}, carries the same rule in its header.

\subsection{What counts as solved}
With $f$ the objective reported by \nirnay{} and $f^\star$ the reference objective,
\begin{equation}\label{eq:relerr}
  \operatorname{err} = \frac{|f-f^\star|}{\max(1,|f^\star|)} .
\end{equation}
\begin{itemize}
  \item \textbf{LP:} status \code{optimal} and $\operatorname{err}<10^{-6}$;
  \item \textbf{MILP:} status \code{optimal} (relative gap proved $\le10^{-4}$) and
        $\operatorname{err}\le10^{-4}$;
  \item \textbf{QP:} status \code{optimal}, $\operatorname{err}<10^{-6}$, and the reference itself
        reports \code{Solved} or \code{AlmostSolved}. Problems whose reference failed are listed but
        not scored.
\end{itemize}
The harnesses also record the largest row and bound violations (and, for MILP, the integrality
violation) of the returned $x$ on the original model. A wrong answer therefore cannot pass on the
objective alone. PDLP stops at a relative KKT tolerance of $10^{-4}$, and its objective errors are reported
against several thresholds (\cref{sec:pdlp-results}).

\subsection{Optimality conditions after postsolve}
\file{tools/dual\_check.py} tests the conditions \eqref{eq:kkt-lp} directly on the original model: primal
feasibility, and the sign of every $y_i$ and $z_j$ against the activity of its row or bound.
This checks presolve, postsolve, unscaling and the simplex duals together
(\cref{sec:presolve-bug}).
\ifnum\haveDualCheck=1 It currently passes on \nDualOk{} of \nDualTot{} Netlib problems with no
failure; the other \nDualStatus{} stop at its \SI{60}{s} limit.\fi

\subsection{Regression tests and file round trips}
\file{tests/} holds \nTestFunctions{} test functions (many parametrised over several instances). They
check each engine against published optimal values: Koch's final Netlib results~\cite{koch2004}, the
MIPLIB~3 solution values and the Maros--Mészáros repository~\cite{maros1999}. They also check the
kernels (the $LDL\T$ on a random quasi-definite matrix, propagation on a hand-made example), and
that every Gomory cut generated on \code{p0033} keeps the known optimal integer point feasible.
Further tests cover PDLP at three tolerances, on infeasible and unbounded problems, with limits, and
on the GPU. The MPS reader and writer are checked by round trips: read, write and read again gives
identical arrays on the benchmark files, and HiGHS reads the original and the written file into the
same model. The case-study tests re-solve every case file with HiGHS against its stored reference.

\subsection{How speed is measured}
Times are wall-clock seconds on the development laptop (Intel Core i5-11400H, 6 cores / 12 threads,
16\,GB RAM, NVIDIA RTX 3050 Laptop GPU with 4\,GB, Windows~11), taken while other work was running.
Differences below a factor of two on sub-second runs are noise. Two aggregate measures are used:
\begin{itemize}
  \item the \textbf{shifted geometric mean} $\operatorname{sgm}_s(t)=\exp\bigl(\frac1N\sum_i\log(t_i+s)\bigr)-s$,
        with shift $s=\shiftA$\,s (and $s=\shiftB$\,s as a check). The shift keeps sub-millisecond
        instances from dominating a ratio of means. It also compresses ratios on small instances,
        so we additionally report the plain geometric mean and the median of per-instance time ratios;
  \item \textbf{performance profiles}~\cite{dolan2002}: for solver $s$ on problem $p$ let
        $r_{p,s}=t_{p,s}/\min_{s'}t_{p,s'}$ ($\infty$ if $s$ failed). The profile
        $\rho_s(\tau)=\frac1N|\{p: r_{p,s}\le\tau\}|$ is the fraction of problems $s$ solves within a
        factor $\tau$ of the fastest. $\rho_s(1)$ is how often $s$ is fastest, and $\rho_s(\infty)$ how
        often it solves at all.
\end{itemize}
